\chapter{扩散模型背后的数学：伊藤微积分与吉尔萨诺夫定理}\label{app:Ito}


基于分数的扩散模型建立在 SDE 上。漂移项描述确定性的运动趋势，布朗运动项引入随机扰动。布朗运动的路径几乎必然处处不可微，通常的链式法则无法直接使用。下面介绍两个处理这种随机性所需的工具。
\begin{itemize}
\item \textbfs{伊藤公式}是随机过程的链式法则。给定满足 SDE 的过程 $\rvx_t$，它描述函数 $\rvh(\rvx_t,t)$ 的演化。用它可以推导福克尔-普朗克方程、矩的演化方程、伊藤乘积法则，以及分数模型训练中的一些恒等式。
\item \textbfs{吉尔萨诺夫定理}描述如何对\emph{路径概率}作测度变换。扩散系数保持不变时，它刻画了改变漂移会怎样改变路径的似然。这将似然加权分数匹配与路径空间 KL 散度联系起来，并通过变分界限与数据似然联系起来。
\end{itemize}

使用这些工具，标准扩散模型推导（福克尔-普朗克方程、反向时间 SDE、训练目标和似然关系）可以清晰地进行，无需依赖非严格的论证。

\chapterlearning{
  \item 理解经典链式法则为何不适用于 SDE，以及伊藤公式如何给出随机过程的链式法则。

  \item 用伊藤公式推导福克尔-普朗克方程及线性 SDE 的解。

  \item 将吉尔萨诺夫定理理解为一种比较随机路径在漂移不同但扩散相同时的概率的方法；

  \item 区分每次匹配分布与匹配完整随机路径，并了解为什么似然加权分数匹配与路径空间 KL 散度和数据似然相关。
}{
本附录介绍扩散模型常用的两种数学工具。伊藤微积分给出了随机微分的运算规则，吉尔萨诺夫定理刻画了漂移项的变化如何影响路径的概率分布。扩散模型的许多推导都用到了这两个工具，有时会省略具体步骤。
}


\section{伊藤公式：随机过程的链式法则}
标准微积分并不直接应用于随机过程，因为维纳过程在经典意义上是不可微的。相反，我们使用伊藤微积分，它提供了处理随机积分的规则。

\subsection{动机：为什么我们需要特殊的链式法则？}
考虑确定性时变函数 $\rvy_t$ 随着时间的推移平稳发展 $t$ （例如，ODE）。如果我们有一个函数 $\rvh(\rvy_t, t)$，通常的链式法则告诉我们：
\[
\frac{\diff \rvh}{\diff t} = \frac{\partial \rvh}{\partial t}  + \nabla_{\rvy}\rvh  \frac{\diff \rvy_t}{\diff t} .
\]
在这里， $\nabla_{\rvy}\rvh$ 是 $\rvh$ 的雅可比矩阵。这非常适合确定性路径 $\rvy_t$.

\begin{question}
那么如果 $\rvx_t$ 是一个随机过程，比如由 \Cref{eq:sde} 那样的 SDE 驱动
\[
\diff \rvx_t = \rvf(\rvx_t, t)\diff t + g(t) \diff \rvw_t
\]
情形又会怎样？过程 $\rvh(\rvx_t, t)$ 满足哪一条 SDE？
\end{question}

\paragraph{为什么普通链式法则失败？}
天真地应用经典链式法则会产生
\[
\diff \rvh = \frac{\partial \rvh}{\partial t} \diff t + \nabla_{\rvy}\rvh  \cdot \diff \rvx_t.
\]
但是，布朗增量按比例缩放
$\mathcal{O}(\sqrt{\diff t})$，而它们的二次变分遵循伊藤微积分的规则
\[
(\diff \rvw_t)^2 = \diff t.
\]
因此，二阶项 $\diff \rvw_t$ \textit{不} 在随机微积分中消失，与经典微积分不同 $(\diff t)^2$ 项可以忽略不计。

\exm{简单示例--$h(x_t) = x_t^2$}{
为了直观，让我们考虑简单的实值函数 $h(x_t) = x_t^2 \in \mathbb{R}$ 其中随机变量 $x_t \in \mathbb{R}$ 满足
\[
\diff x_t = \sigma \diff w_t, \quad
x_0=0,
\]
具有常数 $\sigma>0$。如果我们尝试经典的链式法则，
\[
\diff h = 2 x_t \diff x_t = 2 x_t \sigma \diff w_t.
\]
如果这是真的，则期望 $h(x_t)$ 在时间上是恒定的，因为
\[
\mathbb{E}[\diff h ] =  2 \sigma \mathbb{E}[x_t  \diff w_t]= 2 \sigma \mathbb{E}[x_t]  \underbrace{\mathbb{E}[\diff w_t]}_{=0}=0.
\]
但是我们从经典布朗运动属性中知道（参见 \Cref{eq:wiener-gaussian}）
\[
\mathbb{E}[x_t^2] = \sigma^2 t,
\]
随时间线性增长。所以普通的链式法则遗漏了一个重要的术语。
}


\subsection{由泰勒展开推导一维伊藤公式}\label{subsec:1D-Ito}

\paragraph{通过泰勒展开式确定性链式法则。}
为了理解为什么经典链式法则对于 SDE 定义的随机过程失败，我们首先使用泰勒展开在确定性设置中重新审视它。我们考虑标量情况： $y_t \in \mathbb{R}$ 和 $h(\cdot,\cdot)\in \mathbb{R}$。正式治疗 $\diff y_t = y_{t+\diff t} - y_t$，与 $\diff t \approx 0$，我们扩展：
\begin{align*}
&h(y_{t+\diff t}, t+\diff t)
- h(y_t, t) \\
= ~&\frac{\partial h}{\partial t} \diff t
+ \frac{\partial h}{\partial y} \diff y_t  + \frac{1}{2} \left(\textcolor{gray}{\frac{\partial^2 h}{\partial y^2}  (\diff y_t)^2}
+ 2 \textcolor{gray}{\frac{\partial^2 h}{\partial t \partial y} \diff t \diff y_t}
+  \textcolor{gray}{\frac{\partial^2 h}{\partial t^2} (\diff t)^2}
\right)
+ \mathcal{O}(\diff t^3),
\end{align*}
在这里， $\diff t \diff y_t = \left(\frac{\diff y_t}{\diff t}\right)(\diff t)^2 = \mathcal{O}(\diff t^2)$，类似地 $(\diff t)^2 = \mathcal{O}(\diff t^2)$。因此，所有灰色部分都可以忽略，完整的微分为：
\[
\diff h
= \frac{\partial h}{\partial t} \diff t
+ \frac{\partial h}{\partial y} \diff y_t
+ \mathcal{O}(\diff t^2).
\]


\paragraph{伊藤 的随机泰勒展开式公式。}
现在考虑随机过程 $x_t \in \mathbb{R}$ 受上海证券交易所管辖：
\[
\diff x_t = f(x_t, t) \diff t + g(t) \diff w_t,
\]
其中，$w_t$ 是标准布朗运动。我们要求标量函数 $h(x_t, t)$ 的微分。对 $\diff x_t$ 做随机泰勒展开并保留到二阶~\citep{kloeden1992stochastic}，得到
\begin{align*}
&h(x_{t+\diff t}, t+\diff t)
- h(x_t, t) \\
=~&\frac{\partial h}{\partial t} \diff t
+ \frac{\partial h}{\partial x} \diff x_t
+ \frac{1}{2} \left( \frac{\partial^2 h}{\partial x^2} (\diff x_t)^2
+ 2\textcolor{gray}{{\frac{\partial^2 h}{\partial t \partial x} \diff t \diff x_t}}
+ \textcolor{gray}{{ \frac{\partial^2 h}{\partial t^2} (\diff t)^2}} \right)
+ \cdots
\end{align*}
\subparagraph{可忽略的交叉项。}
通过布朗运动的标度性质（\Cref{eq:wiener-gaussian}),
\[
\diff w_t = \mathcal{O}(\sqrt{\diff t}) \quad \Rightarrow \quad \diff t \cdot \diff w_t = \mathcal{O}((\diff t)^{3/2}).
\]
因此，
\begin{align*}
\diff t \cdot \diff x_t
&= \diff t \left(f \diff t + g \diff w_t\right)
= f (\diff t)^2 + g \cdot \diff t \cdot \diff w_t
= \mathcal{O}((\diff t)^{3/2}).
\end{align*}
所以，灰色项可以忽略不计： $ \mathcal{O}((\diff t)^{3/2}) $ 或更小。
\subparagraph{二阶项 $(\diff x_t)^2$.}
使用 SDE 扩展：
\begin{align*}
(\diff x_t)^2
&= \left( f \diff t + g \diff w_t \right)^2 \\
&= f^2 (\diff t)^2 + 2fg \diff t \diff w_t + g^2 (\diff w_t)^2 \\
&= g^2(t) \diff t,
\end{align*}
最后一个等式使用 伊藤 微分规则
$(\diff t)^2=0$, $\diff t\,\diff w_t=0$, 和
$(\diff w_t)^2=\diff t$，最后一个恒等式在二次变分意义上被理解。

结合各项，我们得到微分：
\[
\diff h(x_t, t)
= \frac{\partial h}{\partial t} \diff t
+ \frac{\partial h}{\partial x} \diff x_t
+ \frac{1}{2} \frac{\partial^2 h}{\partial x^2} g^2(t) \diff t.
\]
替换 $ \diff x_t = f(x_t, t) \diff t + g(t) \diff w_t $ 产量：
\[
\diff h(x_t, t)
= \left( \frac{\partial h}{\partial t} + f \frac{\partial h}{\partial x}
+ \frac{1}{2} g^2 \frac{\partial^2 h}{\partial x^2} \right) \diff t
+ g \frac{\partial h}{\partial x} \diff w_t.
\]
这是 1D 版本 \emph{伊藤公式}.


\exm{简单示例--$h(x_t) = x_t^2$}{
我们回顾一下这个简单的例子： $ h(x_t) = x_t^2 $，其中随机过程 $ x_t \in \mathbb{R} $ 满足
\[
\diff x_t = \sigma \diff w_t,
\]
具有常数 $ \sigma > 0 $。把伊藤公式应用于 $ h(x_t) = x_t^2 $，我们得到：
\[
\diff h(x_t) = \diff (x_t^2) = 2 x_t \diff x_t + \sigma^2 \diff t.
\]
替换 $ \diff x_t = \sigma \diff w_t $，这变成：
\[
\diff(x_t^2) = 2 x_t \sigma \diff w_t + \sigma^2 \diff t.
\]
}

\subsection{伊藤公式：SDE 的链式法则}
\label{subsec:ito-summary}
下面先汇总上文推导的一维伊藤公式，再给出可用类似方法推出的多维公式。最后用它推导\emph{伊藤乘积法则}，计算随机过程 $\rvx_t$ 和 $\rvy_t$ 的内积微分 $\diff(\rvx_t^\top \rvy_t)$。

\paragraph{$1$D 伊藤公式。}
让 $x_t \in \mathbb{R}$ 是满足 SDE 的随机过程：
\[
\diff x_t = f(x_t, t) \diff t + g(t) \diff w_t.
\]
对于标量函数 $h \colon \mathbb{R} \times [0, T] \to \mathbb{R}$，过程 $h(x_t, t)$ 满足：
\begin{mdframed}
\[
\diff h(x_t, t)
= \left( \frac{\partial h}{\partial t} + f \frac{\partial h}{\partial x}
+ \frac{1}{2} g^2 \frac{\partial^2 h}{\partial x^2} \right) \diff t
+ g \frac{\partial h}{\partial x} \diff w_t.
\]
\end{mdframed}


\paragraph{标量值函数的多维伊藤公式。}
让 $\rvx_t \in \mathbb{R}^D$ 满足 SDE：
\[
\diff \rvx_t = \rvf(\rvx_t, t) \diff t + g(t) \diff \rvw_t,
\]
其中，$\rvf \colon \mathbb{R}^D \times [0, T] \to \mathbb{R}^D$ 是漂移项，$g \colon [0, T] \to \mathbb{R}$ 是扩散系数，$\rvw_t \in \mathbb{R}^D$ 是 $D$ 维布朗运动。设 $h \colon \mathbb{R}^D \times [0, T] \to \mathbb{R}$ 为标量函数，则 $h(\rvx_t, t)$ 满足
\begin{mdframed}
\begin{align}\label{eq:Ito-scalar}
    \diff h(\rvx_t, t)
= \left( \frac{\partial h}{\partial t}
+ \nabla_{\rvx} h^\top \rvf
+ \frac{1}{2} g^2(t) \operatorname{Tr} \left( \nabla^2_{\rvx} h \right) \right) \diff t
+ g(t) \nabla_{\rvx} h^\top \diff \rvw_t,
\end{align}
\end{mdframed}
其中，$\nabla_{\rvx} h \in \mathbb{R}^D$ 和 $\nabla^2_{\rvx} h \in \mathbb{R}^{D \times D}$ 分别是 $h$ 关于 $\rvx$ 的梯度和 黑塞矩阵。


\exm{伊藤乘积法则}{
让 $\rvx_t, \rvy_t \in \mathbb{R}^D$ 是由 SDE 控制的向量值随机过程：
\[
\begin{aligned}
\diff \rvx_t &= \rva(\rvx_t, t) \diff t + b(t) \diff \rvw_t, \\
\diff \rvy_t &= \rvc(\rvy_t, t) \diff t + d(t) \diff \rvw_t,
\end{aligned}
\]
其中，$\rva, \rvc: \mathbb{R}^D \times [0, T] \to \mathbb{R}^D$ 是向量场，$b(t), d(t) \in \mathbb{R}$ 是标量函数，$\rvw_t \in \mathbb{R}^D$ 是标准 $D$ 维布朗运动。下面推导标量过程所满足的 SDE
\[
z(t) := \rvx_t^\top \rvy_t.
\]

把多维伊藤公式应用于双线性函数 $h(\rvx, \rvy) := \rvx^\top \rvy$，我们得到：
\[
\diff(\rvx^\top \rvy) = (\diff \rvx)^\top \rvy + \rvx^\top \diff \rvy + \Tr\left[ \diff \rvx \cdot (\diff \rvy)^\top \right].
\]

伊藤 修正项计算如下：
\begin{align*}
    \diff \rvx \cdot (\diff \rvy)^\top &= b(t) \diff \rvw_t \cdot \left[ d(t) \diff \rvw_t \right]^\top \\&= b(t) d(t) \diff \rvw_t \cdot \diff \rvw_t^\top \\&= b(t)d(t) \diff t \cdot \rmI_D.
\end{align*}
因此，
\[
\Tr\left[ \diff \rvx \cdot (\diff \rvy)^\top \right] = b(t)d(t) \Tr(\rmI_D) \diff t = D  b(t)d(t) \diff t.
\]

将所有内容放在一起，生成的 SDE 为：
\begin{align}\label{eq:ito-product}
\begin{aligned}
        \diff(\rvx^\top \rvy) &= (\diff \rvx)^\top \rvy + \rvx^\top \diff \rvy + D  b(t)d(t) \diff t
    % \\ &=\rva^\top \rvy \diff t + b(t) \rvy^\top \diff \rvw_t + \rvx^\top \rvc \diff t + d(t) \rvx^\top \diff \rvw_t + D  b(t)d(t) \diff t.
\end{aligned}
\end{align}
}


\subsection{伊藤公式的应用：福克尔-普朗克方程的推导}\label{subsec:rigorous-proof-fpe}
从 \Cref{eq:Ito-scalar} 中的伊藤公式可以推出福克尔-普朗克方程。这一偏微分方程描述了 \Cref{eq:sde} 所定义的 $ D $ 维扩散过程的概率密度 $ p_t(\rvx) $ 随时间的演化。

\paragraph{第 1 步：应用伊藤（It\^o）公式。}
让 $\phi(\rvx, t)$ 是一个流畅的检验函数 $\phi \colon \mathbb{R}^D \times [0, T] \to \mathbb{R}$。靠伊藤公式：
\[
\diff \phi(\rvx_t, t) = \left( \frac{\partial \phi}{\partial t} + \nabla_{\rvx} \phi^\top \rvf(\rvx_t, t) + \frac{1}{2} g^2(t) \Tr[\nabla_{\rvx}^2 \phi] \right) \diff t + g(t) \nabla_{\rvx} \phi^\top \diff \rvw_t.
\]
\paragraph{第 2 步：获取期望。}
获取期望并使用伊藤 在标准可积条件下积分的期望为零：
\[
\mathbb{E}[\diff \phi(\rvx_t, t)] = \mathbb{E} \left[ \left( \frac{\partial \phi}{\partial t} + \nabla_{\rvx} \phi^\top \rvf(\rvx_t, t) + \frac{1}{2} g^2(t) \Tr[\nabla_{\rvx}^2 \phi] \right) \diff t \right].
\]
\paragraph{第 3 步：通过密度表达期望。}
这个期望可以写成：
\[
\mathbb{E}[\diff \phi(\rvx_t, t)] = \int \left( \frac{\partial \phi}{\partial t} + \nabla_{\rvx} \phi^\top \rvf(\rvx, t) + \frac{1}{2} g^2(t) \Tr[\nabla_{\rvx}^2 \phi] \right) p_t(\rvx) \diff \rvx \diff t.
\]
\paragraph{第 4 步：分部分集成。}
使用分部积分（发散定理） $\mathbb{R}^D$:
\begin{align*}
    \int \nabla_{\rvx} \phi^\top \rvf p_t \diff \rvx &= -\int \phi \nabla_{\rvx} \cdot (\rvf p_t) \diff \rvx, \\
    \int \Tr[\nabla_{\rvx}^2 \phi] p_t  \diff \rvx &= \int \phi \Delta p_t \diff \rvx.
\end{align*}
\paragraph{第 5 步：替换并重新排列。}
请注意
\[
\mathbb{E}[\phi(\rvx_t,t)] = \int \phi(\rvx,t) p_t(\rvx) \diff\rvx.
\]
对时间微分，我们有
\[
\frac{\diff}{\diff t}\mathbb{E}[\phi(\rvx_t,t)]
= \frac{\diff}{\diff t}\int \phi(\rvx,t) p_t(\rvx) \diff\rvx
= \int\Bigl(\partial_t\phi(\rvx,t) p_t(\rvx) + \phi(\rvx,t) \partial_t p_t(\rvx)\Bigr)\diff\rvx.
\]
因此，
\[
\mathbb{E}[\diff \phi(\rvx_t,t)]
= \frac{\diff}{\diff t}\mathbb{E}[\phi(\rvx_t,t)] \diff t
= \int\Bigl(\partial_t\phi(\rvx,t) p_t(\rvx) + \phi(\rvx,t) \partial_t p_t(\rvx)\Bigr)\diff\rvx \diff t.
\]
将此表达式与 Step 等同~3、条款
$\int \partial_t\phi\,p_t\,\diff\rvx\,\diff t$ 取消。申请步骤~4 剩下的空间项给出
\[
\int
\phi(\rvx,t)
\left[
\frac{\partial p_t}{\partial t}
+
\nabla_{\rvx}\cdot\bigl(\rvf(\rvx,t)p_t(\rvx)\bigr)
-
\frac{1}{2}g^2(t)\Delta p_t(\rvx)
\right]
\diff\rvx
=
0.
\]
因为这适用于每个平滑测试函数 $\phi$，我们得到福克尔-普朗克方程
\[
\frac{\partial p_t(\rvx)}{\partial t}
=
-\nabla_{\rvx}\cdot\bigl(\rvf(\rvx,t)p_t(\rvx)\bigr)
+
\frac{1}{2}g^2(t)\Delta p_t(\rvx).
\]

\subsection{伊藤公式应用：线性 SDE 的闭式解}\label{subsec:rigorous-proof-linear-sde}
本小节演示如何通过使用积分因子（类似于 ODE 情况）和伊藤（It\^o）公式来获得线性 SDE 的闭式解。该方法反映了求解线性常微分方程的经典技术，但适应了随机设置。我们考虑以下形式的线性 SDE
\begin{align}\label{eq:app-linear-sde}
    \mathrm{d}\rvx_t = f(t)\rvx_t \mathrm{d}t + g(t) \mathrm{d}\rvw_t,
\end{align}
其中，$f(t)$ 和 $g(t)$ 是确定性函数，$\rvw_t$ 是标准维纳过程。
\paragraph{线性 SDE 的闭式解。} 我们推导出线性 SDE 的显式解 \Cref{eq:app-linear-sde} 采用因子积分法。这种类型的正向 SDE 通常出现在扩散模型中（参见 \Cref{sec:dm-today}).


\subparagraph{步骤 1：定义积分因子。} 让
\[
\Psi(t) := \exp\left( -\int_0^t f(s) \mathrm{d}s \right),
\quad \text{并令} \quad
\rvy_t := \Psi(t)\rvx_t.
\]
\subparagraph{第 2 步：应用伊藤（It\^o）公式。}
把伊藤公式应用于函数 $\rvh(\rvx, t) := \Psi(t)\rvx$。这是伊藤乘积法则的一个特例，见 \Cref{eq:ito-product}。自从 $\Psi(t)$ 是确定性的，没有交叉变分项，公式简化为：
\begin{align*}
    \mathrm{d}\rvy_t
    &= \mathrm{d}[\Psi(t)\rvx_t] \\
    &= \Psi'(t)\rvx_t \mathrm{d}t + \Psi(t) \mathrm{d}\rvx_t \\
    &= -f(t)\Psi(t)\rvx_t \mathrm{d}t + \Psi(t)\left[f(t)\rvx_t \mathrm{d}t + g(t) \mathrm{d}\rvw_t\right] \\
    &= \Psi(t)g(t) \mathrm{d}\rvw_t.
\end{align*}
因此，
\[
\rvy_t = \rvy_0 + \int_0^t \Psi(s)g(s) \mathrm{d}\rvw(s) = \rvx_0 + \int_0^t \Psi(s)g(s) \mathrm{d}\rvw(s),
\]
自 $\Psi(0) = 1$.
\subparagraph{第 3 步：求解 $\rvx_t$.}
使用 $\rvx_t = \Psi(t)^{-1}\rvy_t$，我们得到
\begin{align}\label{eq:lienar-sde-solution}
    \rvx_t &=  e^{ \int_0^t f(s)\mathrm{d}s}
    \left[ \rvx_0 + \int_0^t e^{ -\int_0^s f(r) \mathrm{d}r } g(s) \mathrm{d}\rvw(s) \right].
\end{align}
这为向量值 SDE 提供了显式解。下面，我们演示了两种计算解析形式的替代方法 $p_t(\rvx_t|\rvx_0)$.\paragraph{ 闭式解的分析。} \Cref{eq:lienar-sde-solution} 再次确认 $p_t(\rvx_t|\rvx_0)$ 是高斯的。要看到这一点，定义
\[
\phi(s) := e^{-\int_0^s f(u) \diff u} g(s),
\]
这是时间的确定性矩阵值函数（假设 $f(u)$ 和 $g(s)$ 是确定性的）。 伊藤 积分 $\int_0^t \phi(s) \diff \rvw_s$ 是一个零均值高斯随机变量，因为它是确定性函数关于布朗运动的随机积分。因此，$\rvx_t|\rvx_0$ 是高斯随机变量的仿射变换，因此本身也是高斯的。其分布完全由其条件均值和协方差来表征。

我们将条件均值和协方差（给定初始条件 $\rvx_0$）定义为
\[
\rvm(t) := \mathbb{E}[\rvx_t | \rvx_0], \quad \rmP(t) := \mathbb{E}[(\rvx_t - \rvm(t))(\rvx_t - \rvm(t))^\top|\rvx_0].
\]

\subparagraph{意思是.}
利用期望的线性度以及确定性函数的 伊藤 积分具有零均值这一事实：
\begin{align*}
    \rvm(t) = \mathbb{E}\left[ e^{\int_0^t f(s) \diff s} \left( \rvx_0 + \int_0^t \phi(s) \diff \rvw_s \right) \bigg| \rvx_0 \right] = e^{\int_0^t f(s) \diff s} \rvx_0.
\end{align*}

\subparagraph{协方差.}
让$\rvz_t := \int_0^t \phi(s) \diff \rvw_s$。那么$\rvx_t - \rvm(t) = A(t) \rvz_t$，所以
\[
\rmP(t) = e^{2\int_0^t f(s) \diff s} \mathbb{E}[\rvz_t \rvz_t^\top].
\]
通过 伊藤等距性\footnote{伊藤等距性将随机积分与期望中的标准积分联系起来；我们省略证明，因为它需要 伊藤 微积分的完整机制。对于过程 $\bm{\psi} : [0,T] \to \mathbb{R}^{D \times D}$，伊藤等距性状态
\[
\mathbb{E}\left[\left\| \int_0^T \bm{\psi}(t)   \mathrm{d}\rvw_t \right\|^2 \right]
= \mathbb{E}\left[\int_0^T \|\bm{\psi}(t)\|_F^2   \mathrm{d}t \right],
\]
其中 $\|\bm{\psi}(t)\|_F^2 = \sum_{i,j=1}^D |\psi_{ij}(t)|^2$ 是弗罗贝尼乌斯范数。对于 $\bm{\psi}(t) \in \mathbb{R}^D$（向量），积分是标量，等距简化为
\[
\mathbb{E}\left[\left(\int_0^T \bm{\psi}(t)   \mathrm{d}\rvw_t\right)^2\right] = \mathbb{E}\left[\int_0^T \|\bm{\psi}(t)\|^2   \mathrm{d}t \right].
\]
}，
\[
\mathbb{E}[\rvz_t \rvz_t^\top] = \left( \int_0^t \phi^2(s) \diff s \right) \rmI_D,
\]
因此，
\[
\rmP(t) = e^{2\int_0^t f(s) \diff s} \left( \int_0^t \left( e^{-\int_0^s f(u) \diff u} g(s) \right)^2 \diff s \right) \rmI_D.
\]
这表明条件协方差是各向同性的。

\paragraph{\Cref{eq:mean-var-evolution}.} 中均值和方差 ODE 的推导
或者，我们可以直接从线性 SDE \Cref{eq:app-linear-sde} 导出矩演化方程。


\subparagraph{均值进化.}
取 SDE 两边的条件期望并使用线性：
\[
\frac{\diff \rvm(t)}{\diff t} = \mathbb{E}[f(t) \rvx_t|\rvx_0] = f(t) \mathbb{E}[\rvx_t|\rvx_0] = f(t) \rvm(t).
\]

\vspace{0.3cm}

\subparagraph{协方差演化.}
定义中心化过程 $ \tilde{\rvx}_t := \rvx_t - \rvm(t) $。应用伊藤（It\^o）乘积法则（参见\Cref{eq:ito-product}）：
\[
\diff \left( \tilde{\rvx}_t \tilde{\rvx}_t^\top \right)
= \diff \tilde{\rvx}_t \cdot \tilde{\rvx}_t^\top
+ \tilde{\rvx}_t \cdot \diff \tilde{\rvx}_t^\top
+ \diff \tilde{\rvx}_t \cdot \diff \tilde{\rvx}_t^\top.
\]

根据 SDE，我们计算：
\[
\diff \tilde{\rvx}_t = \diff \rvx_t - \diff \rvm(t) = f(t) \tilde{\rvx}_t  \diff t + g(t)  \diff \rvw_t.
\]

代入乘积法则并取期望：
\begin{align*}
    \frac{\diff \rmP(t)}{\diff t}
    &= \mathbb{E}[f(t) \tilde{\rvx}_t \tilde{\rvx}_t^\top + \tilde{\rvx}_t \tilde{\rvx}_t^\top f(t) + g^2(t) \rmI_D] \\
    &= 2 f(t) \rmP(t) + g^2(t) \rmI_D.
\end{align*}
因此，我们恢复了\Cref{eq:mean-var-evolution}中的矩演化方程。


\section{测度变换：扩散模型中的吉尔萨诺夫定理}
扩散模型利用 SDE 将简单的噪声分布转化为数据分布。保持扩散系数不变，只改变确定性的漂移项，就可以改变过程的路径分布。吉尔萨诺夫定理说明了如何比较这些分布。

\paragraph{两种漂移下的路径似然。}
设观测到从 $t=0$ 到 $t=T$ 的一条连续轨迹 $\rvx_{0:T}:=\{ \rvx_t|t \in [0,T] \}$。我们希望知道两个 SDE 对这条轨迹给出的相对似然。
\begin{question}
    如果两种模型的漂移不同，它们对同一条观测轨迹给出的似然之比是多少？
\end{question}
比较的两个模型具有相同的扩散系数，并对 $\rvx_0$ 采用相同的初始分布。它们分别用漂移 $\rvf$ 和 $\tilde{\rvf}$ 定义路径测度。定理给出这两个测度之间的关系。

\paragraph{设定。}
记 $P_{\rvf}$ 和 $P_{\tilde{\rvf}}$ 为以下两个 SDE 对应的路径测度。
\begin{align*}
    \diff\rvx_t
    &= \rvf(\rvx_t,t)\diff t + g(t)\diff\rvw_t,
    \\
    \diff\rvx_t
    &= \tilde{\rvf}(\rvx_t,t)\diff t + g(t)\diff\rvw_t.
\end{align*}
为简化讨论，假设在所考察的区间内 $g(t)>0$，且满足吉尔萨诺夫定理的标准可积性条件。以 $P_{\tilde{\rvf}}$ 为参考路径测度，记该测度下的布朗运动为 $\rvw_t$，并定义漂移之差
\[
\bm{\delta}_t
:=
\rvf(\rvx_t,t)-\tilde{\rvf}(\rvx_t,t).
\]

\paragraph{路径似然比。}
吉尔萨诺夫定理给出两个路径测度之间的拉东-尼科迪姆导数
\begin{mdframed}
\[
\frac{\diff P_{\rvf}}{\diff P_{\tilde{\rvf}}}(\rvx_{0:T})
=
\exp\left(
\int_0^T
\bm{\delta}_t^\top g(t)^{-1}\diff\rvw_t
-
\frac{1}{2}
\int_0^T
\big\|g(t)^{-1}\bm{\delta}_t\big\|^2
\diff t
\right).
\]
\end{mdframed}
路径似然比的指数由两个积分组成，第一项是伊藤积分，第二项是通常的时间积分。这个比值把随机动力学在路径空间中的差异与基于似然的训练目标联系起来。

这个导数的作用类似于变量变换公式中的雅可比行列式。在扩散系数固定而漂移变化时，拉东-尼科迪姆导数给出了相对于参考测度重加权的因子，可以据此在两种过程的概率和期望之间转换。

\subsection{吉尔萨诺夫定理作为似然训练和分数匹配之间的桥梁}\label{subsec:girsanov}

在了解吉尔萨诺夫定理如何关联单个概率之后不同漂移假设下的路径，我们现在深入研究其对扩散模型。回想一下扩散模型中的前向 SDE：
\[
\diff \rvx_t = \rvf(\rvx_t, t) \diff t + g(t) \diff \rvw_t,
\]
这会在整个轨迹上产生 \emph{ 路径分布 } $P$
$\rvx_{0:T} := \{\rvx_t\}_{t=0}^T$（即过程的联合律整个时间间隔）。
逆时 SDE，由可学习的分数函数参数化
$\rvs_{\bm{\phi}}(\rvx_t, t)$，由下式给出
\[
\diff \rvx_t = \big[\rvf(\rvx_t, t) - g^2(t) \rvs_{\bm{\phi}}(\rvx_t, t)\big] \diff t
+ g(t) \diff \bar{\rvw}_t,
\]
这又定义了另一个路径分布
$P_{\bm{\phi}}$ 超过轨迹。


\paragraph{扩散模型中的两个概念。}
区分两个随机过程可以在两个级别上进行区分是有用的同意：
\begin{itemize}
    \item \textbfs{概念 1.边缘分布匹配：}
    进程在每个固定位置具有相同的分布 $p_t(\rvx_t)$
    时间 $t$。因此，他们的时间索引“快照”是一致的。

    \item \textbfs{概念 2.关节路径分布匹配：}
    这些过程在整个轨迹上具有相同的概率定律。
    这个更强烈的概念也符合不同时期的状态耦合，包括它们的过渡定律和时间依赖性。
\end{itemize}

匹配全路径法则自动意味着边缘匹配，因为每个$p_t$都是通过将路径定律投影到时间状态上得到的
$t$。反之通常是错误的：两个进程可以具有相同的每次都有边缘，同时在这些边缘之间引入不同的耦合边缘。 SDE 及其相关的 PF-ODE 提供了一个重要的示例：
它们共享相同的边缘密度路径，但通常具有不同的样本轨迹和路径定律。


\paragraph{吉尔萨诺夫定理将分数错误与路径律不匹配联系起来。}
精确的逆时 SDE 满足上述更强的概念。下标准正则条件，当从确切的终端初始化时边缘$p_T$并配备真实分数
$\nabla_{\rvx}\log p_t(\rvx)$，是正向的时间反转扩散，因此再现其完整的时间反转路径定律。

让 $P^{\mathrm{rev}}$ 表示这个精确的反时路径定律，并让
$P_{\bm{\phi}}$ 表示学习到的逆时路径定律用 $\rvs_{\bm{\phi}}$ 替换真实分数并从
$p_{\mathrm{prior}}$。将路径空间 KL 分解为初始路径空间之间的不匹配分布和条件路径定律，并应用吉尔萨诺夫使用相同的方法表达两个过程后，得出后者的定理增加逆时时钟，给出原始扩散时间标签$t$，
\begin{align}\label{eq:girsanov-KL}
\begin{aligned}
\mathcal{D}_{\mathrm{KL}}
\bigl(P^{\mathrm{rev}}\|P_{\bm{\phi}}\bigr)
=
&\;
\mathcal{D}_{\mathrm{KL}}
\bigl(p_T\|p_{\mathrm{prior}}\bigr)
\\
&+
\frac{1}{2}
\int_0^T
g^2(t)
\E_{\rvx_t\sim p_t}
\left[
\left\|
\rvs_{\bm{\phi}}(\rvx_t,t)
-
\nabla_{\rvx}\log p_t(\rvx_t)
\right\|^2
\right]
\diff t .
\end{aligned}
\end{align}
第一项衡量真实终端分布之间的不匹配以及用于初始化生成的先验；当它消失时
$p_T=p_{\mathrm{prior}}$。第二项衡量累积的不匹配在习得的漂移和精确的反向漂移之间。因此，与特定似然权重 $g^2(t)$ 的分数匹配有直接的路径空间解释：最小化加权分数误差最小化精确的反向路径定律和学习到的反向路径定律之间的 KL 散度，
直至固定终端分布项。当去噪分数匹配为使用而不是边缘分数匹配，相应的目标不同由与参数无关的常数。通常不同的时间加权不对应于相同的路径空间 KL。

\paragraph{与数据似然的连接.}
路径空间 KL 还控制生成的和数据边缘。由于轨迹的端点是可测量的预测和 KL 散度在边缘化下不会增加，
\begin{align}\label{eq:girsanov-likelihood}
\mathcal{D}_{\mathrm{KL}}
\bigl(p_{\mathrm{data}}\|p_{\bm{\phi}}\bigr)
\leq
\mathcal{D}_{\mathrm{KL}}
\bigl(P^{\mathrm{rev}}\|P_{\bm{\phi}}\bigr),
\end{align}
其中 $p_{\bm{\phi}}$ 表示数据时间生成的边缘。

将此与 \Cref{eq:girsanov-KL} 相结合表明，似然加权
分数匹配与终端分布不匹配一起提供数据侧 KL 散度的上限，因此也是预期的上限负对数似然达到独立于 $\bm{\phi}$ 的常数
\citep{song2021maximum}。
